% !TEX root = ../demo.tex
% =================================================================
%  templates/figure-tikz.tex
%  Иллюстрации. В эталоне по матанализу картинки НЕ плавающие:
%  они ставятся ровно по месту внутри \begin{center}...\end{center},
%  без \caption и без \label. Стили — из style/tikz-styles.tex.
% =================================================================

\subsection{График функции с осями и сеткой}

\begin{center}
	\begin{tikzpicture}[x=1.1cm, y=3.2cm, >=Stealth]
		\pgfmathsetmacro{\n}{6}
		\pgfmathsetmacro{\nminus}{\n-1}
		\pgfmathsetmacro{\nplus}{\n+1}
		\pgfmathsetmacro{\yn}{ln(\n)/ln(10)}
		\pgfmathsetmacro{\yplus}{ln(\nplus)/ln(10)}
		
		\draw[gridline] (-0.6,-1.05) grid (8.6,1.15);
		\draw[ax] (-0.6,0) -- (8.6,0);
		\draw[ax] (0,-1.05) -- (0,1.15);
		
		\foreach \x/\lab in {1/1,2/2,3/3,\nminus/{$n-1$},\n/{$n$},\nplus/{$n+1$}}{
			\draw[tickmark] (\x,0.03) -- (\x,-0.03);
			\node[below] at (\x,0) {\lab};
		}
		\node[below] at (4.1,0) {$\dots$};
		
		\draw[curve, domain=0.15:8.2, samples=240] plot (\x,{ln(\x)/ln(10)});
		\fill (1,0) circle (1.6pt);
		
		% прямоугольник нижней суммы
		\fill[hatchA] (\nminus,0) rectangle (\n,\yn);
		\draw[curvehl] (\nminus,0) rectangle (\n,\yn);
		\node[violet] at ({(\nminus+\n)/2},{\yn/2}) {$\ln n\cdot 1$};
		
		% прямоугольник верхней суммы
		\fill[hatchB] (\n,0) rectangle (\nplus,\yplus);
		\draw[curveaux] (\n,0) rectangle (\nplus,\yplus);
		
		\draw[-{Stealth[length=2mm]}] (8.1,0.98) -- (7.2,{ln(7.2)/ln(10)});
		\node[right] at (8.1,0.98) {$\ln x$};
	\end{tikzpicture}
\end{center}

Пояснение к рисунку идёт обычным абзацем сразу после него --- подписи
\verb|\caption| в этом конспекте не используются.

\subsection{Векторы и ортогональная проекция}

\begin{center}
	\begin{tikzpicture}[>=Stealth, scale=1.2]
		\draw[->] (-1,0) -- (4.4,0) node[right] {$\operatorname{span}\{e_k\}_{k=1}^{N}$};
		\coordinate (O) at (0,0);
		\coordinate (P) at (3,0);
		\coordinate (X) at (3,2);
		\draw[vecarr] (O) -- (X) node[above left] {$x$};
		\draw[vecarr] (O) -- (P) node[below=5pt] {$\sum_{k=1}^{N} c_k e_k$};
		\draw[vecaux] (P) -- (X) node[midway, right] {$x - \sum c_k e_k$};
		\rightangle{(P)}
		\fill (O) circle (1.3pt);
	\end{tikzpicture}
\end{center}

\subsection{Площадь под кривой и разбиение отрезка}

\begin{center}
	\begin{tikzpicture}[x=1.6cm, y=1.6cm, >=Stealth]
		\draw[ax] (-0.3,0) -- (4.6,0) node[right] {$x$};
		\draw[ax] (0,-0.3) -- (0,2.4) node[above] {$y$};
		
		% ступеньки нижней суммы
		\foreach \i in {0,...,3}{
			\pgfmathsetmacro{\xa}{0.5+\i*0.9}
			\pgfmathsetmacro{\xb}{\xa+0.9}
			\pgfmathsetmacro{\h}{0.25*(\xa)*(\xa)+0.3}
			\fill[hatchGray] (\xa,0) rectangle (\xb,\h);
			\draw[thin] (\xa,0) rectangle (\xb,\h);
		}
		
		\draw[curvehl, domain=0.2:4.3, samples=200, smooth]
		plot (\x,{0.25*\x*\x+0.3});
		\node[violet, above left] at (4.3,{0.25*4.3*4.3+0.3}) {$f$};
		
		\foreach \x/\lab in {0.5/{$a$}, 4.1/{$b$}}{
			\draw[tickmark] (\x,0.06) -- (\x,-0.06);
			\node[below] at (\x,0) {\lab};
		}
	\end{tikzpicture}
\end{center}

\subsection{Окрестность точки на прямой}

\begin{center}
	\begin{tikzpicture}[x=1.4cm, >=Stealth]
		\draw[ax] (-0.4,0) -- (5.2,0) node[right] {$\RR$};
		\fill[hatchA] (2,-0.12) rectangle (3.6,0.12);
		\draw[curvehl] (2,-0.12) -- (2,0.12);
		\draw[curvehl] (3.6,-0.12) -- (3.6,0.12);
		\node[pt] at (2.8,0) {};
		\node[below=6pt] at (2.8,0) {$x$};
		\node[below=6pt] at (2,0) {$x-\delta$};
		\node[below=6pt] at (3.6,0) {$x+\delta$};
		\node[above=4pt] at (2.8,0.12) {$U_\delta(x)$};
	\end{tikzpicture}
\end{center}

\subsection{Два графика рядом (scope + xshift)}

\begin{center}
	\begin{tikzpicture}[x=0.9cm, y=0.9cm, >=Stealth]
		\begin{scope}
			\draw[ax] (-0.3,0) -- (3.4,0);
			\draw[ax] (0,-1.4) -- (0,1.6);
			\draw[curve, domain=0.15:3.2, samples=200] plot (\x,{sin(deg(3*\x))});
			\node[below] at (1.6,-1.5) {равномерная сходимость};
		\end{scope}
		\begin{scope}[xshift=5cm]
			\draw[ax] (-0.3,0) -- (3.4,0);
			\draw[ax] (0,-1.4) -- (0,1.6);
			\draw[curve, domain=0.15:3.2, samples=200] plot (\x,{sin(deg(9*\x))/(0.4+\x)});
			\node[below] at (1.6,-1.5) {поточечная сходимость};
		\end{scope}
	\end{tikzpicture}
\end{center}

\subsection{Хайлайты в иллюстрациях}

\begin{center}
	\begin{tikzpicture}[x=1.2cm, y=1.2cm, >=Stealth]
		\draw[->] (-0.2,0) -- (3.8,0) node[right] {$x$};
		\draw[->] (0,-0.2) -- (0,2.3) node[above] {$y$};
		\fill[hlGreen, hlarea] (0.6,0) rectangle (2.4,1.32);
		\draw[hlRed, hlcurve, domain=0.15:3.6, samples=120, smooth]
		plot (\x,{0.55*\x});
		\node[hlBlue, hlmark] at (2.4,1.32) {};
		\node[slbl, above left] at (2.4,1.32) {$(x_0,y_0)$};
	\end{tikzpicture}
\end{center}

Цвета палитры доступны в tikz напрямую, а стили \texttt{hlcurve} (толстая
выделенная кривая), \texttt{hlarea} (полупрозрачная заливка) и
\texttt{hlmark} (отмеченная точка) заданы в \texttt{style/highlight.tex}.
